\section{The second Euclidean diffusion of the actual arithmetic observable}
\label{sec:ab-bilaplacian}

The first diffusion coefficient in the preceding calculation has a finite
limit along the escaping inverse-fibre trajectory. We compute the next
Euclidean diffusion without dropping any derivative of that coefficient.
All spatial derivatives in this section refer to the original Cartesian
coordinates $(x,y,w)$, with their Euclidean metric. Newman time $t$ remains
distinct from the fluid trajectory time $\tau$.

Keep
\[
 Q=1+xy,\qquad
 \sigma(x,y,w)=\tfrac12\bigl[Q^3w+y^2Q(4+3xy)\bigr],
 \qquad \Psi_t(x,y,w)=Z(t,\sigma(x,y,w)).
\]
Here $Z$ is the actual entire function defined by
\eqref{eq:zt-datum}; its initial datum is the specified completed zeta
function. Its analyticity and all differentiated integral bounds were
proved above. In particular, for each real $t$,
\[
 Z_s(t,0)=Z_{sss}(t,0)=0,\qquad
 Z_{ss}(t,0)=-32M_1(t)<0.
\]
The strict inequality retains the positive original moment
$M_1(t)=\int_0^\infty u^2e^{tu^2}\Phi(u)\,du$.

\subsection{The full fourth-order chain rule}

Write $g=\nabla\sigma$, $H=D^2\sigma$, $N=g\cdot g$, and
$L=\Delta\sigma$, where $\Delta=\partial_x^2+\partial_y^2+\partial_w^2$.
For any entire function $h$ of one variable, ordinary differentiation gives
\begin{align}
 \Delta(h\circ\sigma)&=Nh''(\sigma)+Lh'(\sigma),\label{ab:first}\\
 \Delta^2(h\circ\sigma)&=
 N^2h^{(4)}(\sigma)
 +(4g^{\mathsf T}Hg+2NL)h^{(3)}(\sigma)\notag\\
 &\quad+
 (2\lVert H\rVert_{\mathrm F}^2+4g\cdot\nabla L+L^2)h''(\sigma)
 +(\Delta L)h'(\sigma).\label{ab:second}
\end{align}
Here $\lVert H\rVert_{\mathrm F}^2=\sum_{i,j=1}^3H_{ij}^2$ on the
real locus. To prove every coefficient, differentiate the first formula
by the product rule:
\begin{align*}
 \Delta(Nh'')&=(\Delta N)h''
     +2(g\cdot\nabla N)h^{(3)}+NLh^{(3)}+N^2h^{(4)},\\
 \Delta(Lh')&=(\Delta L)h'
     +2(g\cdot\nabla L)h''+L^2h''+LNh^{(3)}.
\end{align*}
In these lines every derivative of $h$ is evaluated at $\sigma$.
For the remaining two identities, write
$N=\sum_i(\partial_i\sigma)^2$ and differentiate each summand:
\[
 \partial_jN=2\sum_i(\partial_i\sigma)(\partial_j\partial_i\sigma),
 \qquad
 \Delta N=2\sum_{i,j}(\partial_j\partial_i\sigma)^2
             +2\sum_i(\partial_i\sigma)\partial_i\Delta\sigma.
\]
Thus $g\cdot\nabla N=2g^{\mathsf T}Hg$ and
$\Delta N=2\lVert H\rVert_{\mathrm F}^2+2g\cdot\nabla L$.
Substitution proves \eqref{ab:second}, including its second- and
first-derivative terms.

\subsection{All coefficients on the unchanged inverse trajectory}

For $z>0$ set
\[
 \gamma(z)=(z^{-1},-3z/2,13z^2/2),\qquad
 \tau=(1-z^2)/8.
\]
The inverse is $z=\sqrt{1-8\tau}$ on $\tau<1/8$.
This is precisely the original trajectory, parametrized by a positive
coordinate; it has not been rescaled. Direct substitution gives
$Q(\gamma)=-1/2$ and $\sigma(\gamma)=-z^2/8$.
For clarity the original first derivatives and Laplacian before
substitution are
\begin{align*}
 g_1&=\tfrac32yQ^2w+\tfrac12y^3(7+6xy),\\
 g_2&=\tfrac32xQ^2w+\tfrac12(8y+21xy^2+12x^2y^3),\\
 g_3&=\tfrac12Q^3,\\
 L&=3wQ(x^2+y^2)+3y^4+18x^2y^2+21xy+4.
\end{align*}
Differentiating these polynomials, followed by the stated substitution,
gives
\[
 g(\gamma)=
 \begin{pmatrix}-9z^3/32\\-3z/16\\-1/16\end{pmatrix},
 \qquad
 H(\gamma)=
 \begin{pmatrix}
 -27z^4/4&3z^2/16&-9z/16\\
 3z^2/16&13/4&3/(8z)\\
 -9z/16&3/(8z)&0
 \end{pmatrix}.
\]
The full contractions needed in \eqref{ab:second} are
\begin{align*}
 N(\gamma)&=\frac{81z^6+36z^2+4}{1024},&
 L(\gamma)&=\frac{13-27z^4}{4},\\
 \lVert H(\gamma)\rVert_{\mathrm F}^2
 &=\frac{729z^8}{16}+\frac{9z^4}{128}
       +\frac{81z^2}{128}+\frac{169}{16}+\frac9{32z^2},\\
 (g\cdot\nabla L)(\gamma)
 &=\frac{9477z^8}{512}-\frac{405z^4}{64}
       +\frac{27z^2}{128}+\frac{81}{32}+\frac3{32z^2},\\
 (g^{\mathsf T}Hg)(\gamma)
 &=-\frac{2187z^{10}}{4096}+\frac{81z^6}{4096}
       -\frac{81z^4}{4096}+\frac{117z^2}{1024}+\frac9{1024},\\
 (\Delta L)(\gamma)&=-111z^2+\frac{36}{z^2}.
\end{align*}
These are polynomial differentiation and finite matrix products, with
both occurrences of each off-diagonal Hessian entry retained in its
squared Frobenius norm. They also give a direct way to check each term
in the following fully expanded result:
\begin{equation}
 (\Delta^2\Psi_t)(\gamma(z))
 =\sum_{j=1}^4 C_j(z)\,\partial_s^jZ(t,-z^2/8),\label{ab:full}
\end{equation}
where
\begin{align*}
 C_4(z)&=\left(\frac{81z^6+36z^2+4}{1024}\right)^2,\\
 C_3(z)&=-\frac{6561z^{10}}{2048}+\frac{243z^6}{2048}
           -\frac{135z^4}{1024}+\frac{351z^2}{512}+\frac{31}{512},\\
 C_2(z)&=\frac{26973z^8}{128}-\frac{4419z^4}{64}
           +\frac{135z^2}{64}+\frac{669}{16}+\frac{15}{16z^2},\\
 C_1(z)&=-111z^2+\frac{36}{z^2}.
\end{align*}
For example, the coefficient of $z^{-2}$ in $C_2$ is
$2(9/32)+4(3/32)=15/16$; the $L^2$ term has no negative
power. Thus this coefficient includes the differentiated first-order
coefficient as well as the Hessian square.

\begin{proposition}[Two successive Euclidean diffusion limits]
For every fixed real Newman time $t$,
\begin{align}
 \lim_{z\downarrow0}(\Delta\Psi_t)(\gamma(z))
    &=\frac{Z_{ss}(t,0)}{256}
      =-\frac{M_1(t)}8,\label{ab:first-limit}\\
 \lim_{z\downarrow0}z^2(\Delta^2\Psi_t)(\gamma(z))
    &=\frac{15}{16}Z_{ss}(t,0)
      =-30M_1(t)<0.\label{ab:second-limit}
\end{align}
In particular $\Delta^2\Psi_t$ is unbounded on $\mathbb R^3$,
although $\Psi_t$ is real analytic and smooth at every finite
spatial point.
\end{proposition}
\begin{proof}
At this fixed $t$, the entire even Taylor expansion of $Z(t,s)$
and its termwise derivatives give, with $d_j=\partial_s^jZ(t,0)$,
\begin{align*}
 Z_s(t,-z^2/8)&=-d_2z^2/8+O(z^6),\\
 Z_{ss}(t,-z^2/8)&=d_2+O(z^4),\\
 Z_{sss}(t,-z^2/8)&=-d_4z^2/8+O(z^6),\\
 Z_{ssss}(t,-z^2/8)&=d_4+O(z^4).
\end{align*}
The remainders follow from convergence on a fixed closed disk around
zero; they are bounded there by the convergent Taylor tails.
Equation \eqref{ab:first}, $N(\gamma)\to1/256$, and
$L(\gamma)\to13/4$ prove \eqref{ab:first-limit}.
In \eqref{ab:full}, $C_4$ and $C_3$ are bounded as $z\downarrow0$,
$z^2C_2\to15/16$, and $z^2C_1\to36$.
Multiplying that equation by $z^2$ and using the four Taylor estimates
proves \eqref{ab:second-limit}. The last term tends to zero because
$Z_s(t,-z^2/8)\to0$; its unmultiplied finite contribution is
$-36d_2/8+O(z^4)$ and has not been discarded from the exact formula.
The moment identity and positivity give the displayed strict sign.
Consequently, for sufficiently small $z>0$ the bilaplacian is at most
$-15M_1(t)z^{-2}$, proving unboundedness. Entire composition proves
smoothness at each finite point. Finally
$|\gamma(z)|\ge z^{-1}\to\infty$, so this sequence is explicitly
at spatial infinity.
\end{proof}

\subsection{The exact relation to the material diffusion operator}

Retain the original divergence-free polynomial flow $X_\tau$ and its
Jacobian $A_\tau$ from the material-generator construction. On each of
its open flow domains, $X_\tau$ is a smooth diffeomorphism,
$\det A_\tau=1$, and
$K_\tau=A_\tau^{-1}A_\tau^{-\mathsf T}$.
The scalar operator proved there is
$\mathcal L_\tau=\operatorname{div}(K_\tau\nabla)$.
For every smooth physical scalar $v$ on the image domain the full
chain and cofactor identity give
\[
 \mathcal L_\tau(v\circ X_\tau)=(\Delta v)\circ X_\tau.
\]
Apply this same identity a second time, now to the smooth scalar
$\Delta v$. At this fixed $\tau$ it proves the exact intertwining
\[
 \mathcal L_\tau^2(v\circ X_\tau)=(\Delta^2v)\circ X_\tau.
\]
The square is operator composition; it differentiates every spatial
coefficient of $K_\tau$. It is not the square of a coefficient
evaluated at one point.

The already established original flow identity
$F_1(X_\tau(q))=F_1(q)+2\tau$ gives
\[
 \Psi_t(X_\tau(q))=Z(t,\sigma(q)+\tau).
\]
At $q_0=(1,-3/2,13/2)$ and $\tau<1/8$ one has
$X_\tau(q_0)=\gamma(\sqrt{1-8\tau})$.
Therefore \eqref{ab:first-limit} and \eqref{ab:second-limit}
give, without interchanging the two time variables,
\begin{align*}
 \lim_{\tau\uparrow1/8}
 \mathcal L_\tau[Z(t,\sigma+\tau)](q_0)&=-M_1(t)/8,\\
 \lim_{\tau\uparrow1/8}(1-8\tau)
 \mathcal L_\tau^2[Z(t,\sigma+\tau)](q_0)&=-30M_1(t).
\end{align*}
For the physical operator $\nu\Delta$ with the unchanged $\nu>0$,
the respective limits acquire the factors $\nu$ and $\nu^2$.
The scalar in brackets itself has an entire extension in $\tau$.
Its material coefficients nevertheless have the displayed second
diffusion behavior at the end of this incomplete chart.
These identities are an exact relation between the arithmetic observable
and Euclidean viscosity. They do not give a singular velocity at a
finite physical point, admissible initial velocity or forcing, or a
counterexample to the stated Navier--Stokes alternatives.
